\documentclass[11pt,a4paper]{article}
\usepackage[margin=1in]{geometry}
\usepackage[T1]{fontenc}
\usepackage[utf8]{inputenc}
\usepackage[english]{babel}
\usepackage{newtxtext}
\usepackage{newtxmath}
\usepackage[hidelinks]{hyperref}

\title{ARTPS: Operationally Safe Target Prioritization via Depth-Enhanced Hybrid Anomaly Detection}
\author{Poyraz BAYDEM\.{I}R\\Sel\c{c}uk University (Konya, T\"urkiye)}
\date{}

\begin{document}
\maketitle

\begin{abstract}
Planetary rovers operate under strict bandwidth and energy constraints with delayed ground feedback, which motivates reliable onboard prioritization of scientific targets. However, visual-only autonomy is brittle in field imagery: illumination changes, cast shadows, and low-texture regions can trigger false positives, while diversity heuristics can suppress high-value anomalies that appear visually similar to previously selected targets (false negatives).

We present ARTPS, an edge-oriented framework that couples monocular depth estimation with multi-cue anomaly detection and robust localization. ARTPS fuses image cues (texture/edge and photometric indicators) with depth cues (discontinuities and proximity weighting) and optional learned anomaly maps. To avoid fixed, hand-tuned weights, ARTPS uses entropy-weighted dynamic fusion that adapts to field conditions by emphasizing the most confident anomaly components at runtime.

To improve operational safety, ARTPS integrates feature-space clustering for multi-shape management with a soft similarity penalty, and introduces a Priority Buffer mechanism that stores high raw-anomaly candidates down-weighted by diversity policy for second-chance re-evaluation when resources allow. On NASA Mars rover imagery, ARTPS achieves 0.894 AUROC, 0.847 AUPRC, and 0.823 F1, outperforming a PaDiM/PatchCore baseline (0.856 AUROC). In a lightweight configuration (OpenCV preprocessing + fusion + localization, excluding learned depth/AE inference), the core pipeline sustains 28.1 FPS at 256x256 on a development workstation, supporting profile-aware onboard screening and enables science-driven autonomy on constrained onboard hardware without compromising reliability.
\end{abstract}

\end{document}

